% GENERATED FILE. Edit canonical lessons, then run npm run generate:books.
% canonical-source-hash: fnv1a64:413c45f5023fc00d
% canonical-lessons: KA-C490-hedaru, KA-C490-nacu, KA-C490-atanka, KA-C490-mujugara, KA-C490-kalavala

\chapter{Fear and Shyness}
\label{ch:ka-fear-and-shyness}
\hlchaptermodality{490}

\begin{chapteropening}
\textbf{By the end of this chapter:} \emph{I can say to be afraid, to feel shy, anxiety, embarrassment and unease.}

Complete the last lesson of chapter 490: I can say to be afraid, to feel shy, anxiety, embarrassment and unease.
\end{chapteropening}

\section[\texorpdfstring{\kn{ಹೆದರು} (hedaru)}{ಹೆದರು (hedaru)}]{\kn{ಹೆದರು} (hedaru) — to be afraid; to fear}

\label{lesson:KA-C490-hedaru}



\begin{quote}
Before the new one: say the Kannada for chalk, then the Kannada for a ruler.
\end{quote}

\subsection*{You'll want to know: \kn{ಹೆದರು}}
\textbf{\kn{ಹೆದರು}} — \emph{hedaru} — \textquotedblleft{}to be afraid; to fear\textquotedblright{}.

\begin{quote}
\kn{ಮಗು} \kn{ಕತ್ತಲೆಗೆ} \kn{ಹೆದರುತ್ತದೆ}.
\end{quote}

\begin{quote}
\emph{magu kattalege hedaruttade.}
\end{quote}

\begin{quote}
The child is afraid of the dark.
\end{quote}

Past (I …): \kn{ಹೆದರಿದೆ}, \emph{hedaride}.

What you fear takes -\kn{ಗೆ}: \kn{ಕತ್ತಲೆಗೆ} \kn{ಹೆದರು}, to be afraid of the dark. The noun for the feeling is \kn{ಭಯ}.

\begin{cousinweb}
\textbf{\kn{ಹೆದರು}} (\emph{hedaru}) is native Dravidian: DEDR 4401 lists Kannada \emph{bedaṟu} and \emph{hedaṟu}, \textquotedblleft{}to be alarmed, to take fright\textquotedblright{}, with a noun meaning \textquotedblleft{}fright\textquotedblright{} beside them. The b- and h- forms are two shapes of one old verb.
\end{cousinweb}

\begin{grammarlens}[title={What you've built}]
One more word for describing what you lived and felt.
\end{grammarlens}

\subsection*{Your turn}
\begin{itemize}
\raggedright
  \item \emph{Say it:} \emph{hedaru}
  \item \emph{Say it:} \emph{hedaru}, once more
  \item \emph{Say it:} \emph{magu kattalege hedaruttade.} — The child is afraid of the dark.
  \item \emph{Recall:} say the Kannada for chalk, then the Kannada for a ruler, then say \emph{hedaru} again
\end{itemize}

\begin{tcolorbox}[breakable,colback=teal!4,colframe=teal!35!black,title={Before you move on}]
What does \kn{ಹೆದರು} mean? (\textquotedblleft{}To be afraid; to fear\textquotedblright{}.) Say it once more.
\end{tcolorbox}
\section[\texorpdfstring{\kn{ನಾಚು} (nācu)}{ನಾಚು (nācu)}]{\kn{ನಾಚು} (nācu) — to feel shy; to be bashful}

\label{lesson:KA-C490-nacu}



\begin{quote}
Before the new one: say the Kannada for a ruler, then the Kannada for to be afraid.
\end{quote}

\subsection*{You'll want to know: \kn{ನಾಚು}}
\textbf{\kn{ನಾಚು}} — \emph{nācu} — \textquotedblleft{}to feel shy; to be bashful\textquotedblright{}.

\begin{quote}
\kn{ಹೊಸ} \kn{ಅತಿಥಿಯನ್ನು} \kn{ನೋಡಿ} \kn{ಹುಡುಗಿ} \kn{ನಾಚಿದಳು}.
\end{quote}

\begin{quote}
\emph{hosa atithiyannu nōḍi huḍugi nācidaḷu.}
\end{quote}

\begin{quote}
Seeing the new guest, the girl felt shy.
\end{quote}

Past (I …): \kn{ನಾಚಿದೆ}, \emph{nācide}.

The noun beside it is \kn{ನಾಚಿಕೆ}, shyness: the verb is the feeling happening, the noun is the feeling itself.

\begin{cousinweb}
\textbf{\kn{ನಾಚು}} (\emph{nācu}) is native Dravidian: DEDR 3639 gives Kannada \emph{nāṇcu}, \emph{nāñcu} and \emph{nācu}, \textquotedblleft{}to become abashed\textquotedblright{}, beside a noun for bashfulness and modesty.
\end{cousinweb}

\begin{grammarlens}[title={What you've built}]
One more word for describing what you lived and felt.
\end{grammarlens}

\subsection*{Your turn}
\begin{itemize}
\raggedright
  \item \emph{Say it:} \emph{nācu}
  \item \emph{Say it:} \emph{nācu}, once more
  \item \emph{Say it:} \emph{hosa atithiyannu nōḍi huḍugi nācidaḷu.} — Seeing the new guest, the girl felt shy.
  \item \emph{Recall:} say the Kannada for a ruler, then the Kannada for to be afraid, then say \emph{nācu} again
\end{itemize}

\begin{tcolorbox}[breakable,colback=teal!4,colframe=teal!35!black,title={Before you move on}]
What does \kn{ನಾಚು} mean? (\textquotedblleft{}To feel shy; to be bashful\textquotedblright{}.) Say it once more.
\end{tcolorbox}
\section[\texorpdfstring{\kn{ಆತಂಕ} (ātaṅka)}{ಆತಂಕ (ātaṅka)}]{\kn{ಆತಂಕ} (ātaṅka) — anxiety; dread}

\label{lesson:KA-C490-atanka}



\begin{quote}
Before the new one: say the Kannada for to be afraid, then the Kannada for to feel shy.
\end{quote}

\subsection*{You'll want to know: \kn{ಆತಂಕ}}
\textbf{\kn{ಆತಂಕ}} — \emph{ātaṅka} — \textquotedblleft{}anxiety; dread\textquotedblright{}.

\begin{quote}
\kn{ಪರೀಕ್ಷೆಗೆ} \kn{ಮೊದಲು} \kn{ನನಗೆ} \kn{ಆತಂಕ} \kn{ಆಯಿತು}.
\end{quote}

\begin{quote}
\emph{parīkṣege modalu nanage ātaṅka āyitu.}
\end{quote}

\begin{quote}
Before the exam I felt anxious.
\end{quote}

A feeling takes the dative of the person: \kn{ನನಗೆ} \kn{ಆತಂಕ} \kn{ಆಯಿತು}, I grew anxious.

\begin{cousinweb}
\textbf{\kn{ಆತಂಕ}} (\emph{ātaṅka}) is borrowed from Sanskrit \emph{ātaṅka}, which Monier-Williams gives first as a sickness of the body and then as affliction of mind, disquiet and fear (Monier-Williams). Kannada keeps the second sense: the dread before something happens.
\end{cousinweb}

\begin{grammarlens}[title={What you've built}]
One more word for describing what you lived and felt.
\end{grammarlens}

\subsection*{Your turn}
\begin{itemize}
\raggedright
  \item \emph{Say it:} \emph{ātaṅka}
  \item \emph{Say it:} \emph{ātaṅka}, once more
  \item \emph{Say it:} \emph{parīkṣege modalu nanage ātaṅka āyitu.} — Before the exam I felt anxious.
  \item \emph{Recall:} say the Kannada for to be afraid, then the Kannada for to feel shy, then say \emph{ātaṅka} again
\end{itemize}

\begin{tcolorbox}[breakable,colback=teal!4,colframe=teal!35!black,title={Before you move on}]
What does \kn{ಆತಂಕ} mean? (\textquotedblleft{}Anxiety; dread\textquotedblright{}.) Say it once more.
\end{tcolorbox}
\section[\texorpdfstring{\kn{ಮುಜುಗರ} (mujugara)}{ಮುಜುಗರ (mujugara)}]{\kn{ಮುಜುಗರ} (mujugara) — embarrassment; awkwardness}

\label{lesson:KA-C490-mujugara}



\begin{quote}
Before the new one: say the Kannada for to feel shy, then the Kannada for anxiety.
\end{quote}

\subsection*{You'll want to know: \kn{ಮುಜುಗರ}}
\textbf{\kn{ಮುಜುಗರ}} — \emph{mujugara} — \textquotedblleft{}embarrassment; awkwardness\textquotedblright{}.

\begin{quote}
\kn{ಎಲ್ಲರ} \kn{ಮುಂದೆ} \kn{ಹಾಡುವಾಗ} \kn{ನನಗೆ} \kn{ಮುಜುಗರ} \kn{ಆಯಿತು}.
\end{quote}

\begin{quote}
\emph{ellara munde hāḍuvāga nanage mujugara āyitu.}
\end{quote}

\begin{quote}
Singing in front of everyone, I felt embarrassed.
\end{quote}

Sharper than \kn{ನಾಚಿಕೆ}: the discomfort of an awkward moment, often caused by someone else.

\begin{cousinweb}
\textbf{\kn{ಮುಜುಗರ}} (\emph{mujugara}) is an everyday Kannada word; its origin is not settled, and none of the dictionaries we can open lists it.
\end{cousinweb}

\begin{grammarlens}[title={What you've built}]
One more word for describing what you lived and felt.
\end{grammarlens}

\subsection*{Your turn}
\begin{itemize}
\raggedright
  \item \emph{Say it:} \emph{mujugara}
  \item \emph{Say it:} \emph{mujugara}, once more
  \item \emph{Say it:} \emph{ellara munde hāḍuvāga nanage mujugara āyitu.} — Singing in front of everyone, I felt embarrassed.
  \item \emph{Recall:} say the Kannada for to feel shy, then the Kannada for anxiety, then say \emph{mujugara} again
\end{itemize}

\begin{tcolorbox}[breakable,colback=teal!4,colframe=teal!35!black,title={Before you move on}]
What does \kn{ಮುಜುಗರ} mean? (\textquotedblleft{}Embarrassment; awkwardness\textquotedblright{}.) Say it once more.
\end{tcolorbox}
\section[\texorpdfstring{\kn{ಕಳವಳ} (kaḷavaḷa)}{ಕಳವಳ (kaḷavaḷa)}]{\kn{ಕಳವಳ} (kaḷavaḷa) — unease; agitation}

\label{lesson:KA-C490-kalavala}



\begin{quote}
Before the new one: say the Kannada for anxiety, then the Kannada for embarrassment.
\end{quote}

\subsection*{You'll want to know: \kn{ಕಳವಳ}}
\textbf{\kn{ಕಳವಳ}} — \emph{kaḷavaḷa} — \textquotedblleft{}unease; agitation\textquotedblright{}.

\begin{quote}
\kn{ರಾತ್ರಿ} \kn{ಮಗ} \kn{ಬರಲಿಲ್ಲ}, \kn{ಅಮ್ಮನಿಗೆ} \kn{ಕಳವಳ} \kn{ಆಯಿತು}.
\end{quote}

\begin{quote}
\emph{rātri maga baralilla, ammanige kaḷavaḷa āyitu.}
\end{quote}

\begin{quote}
The son had not come by night; his mother grew uneasy.
\end{quote}

A restless worry that will not settle; \kn{ಕಳವಳಗೊಳ್ಳು}, to grow uneasy.

\begin{cousinweb}
\textbf{\kn{ಕಳವಳ}} (\emph{kaḷavaḷa}) is native Dravidian: DEDR 1306 lists Kannada \emph{kaḷakaḷa} and \emph{kaḷavaḷike}, \textquotedblleft{}agitation of the mind, distress\textquotedblright{}, and the verb \emph{kaḷavaḷisu}, \textquotedblleft{}to be agitated\textquotedblright{}.
\end{cousinweb}

\begin{grammarlens}[title={What you've built}]
One more word for describing what you lived and felt.
\end{grammarlens}

\subsection*{Your turn}
\begin{itemize}
\raggedright
  \item \emph{Say it:} \emph{kaḷavaḷa}
  \item \emph{Say it:} \emph{kaḷavaḷa}, once more
  \item \emph{Say it:} \emph{rātri maga baralilla, ammanige kaḷavaḷa āyitu.} — The son had not come by night; his mother grew uneasy.
  \item \emph{Recall:} say the Kannada for anxiety, then the Kannada for embarrassment, then say \emph{kaḷavaḷa} again
\end{itemize}

\begin{tcolorbox}[breakable,colback=teal!4,colframe=teal!35!black,title={Before you move on}]
What does \kn{ಕಳವಳ} mean? (\textquotedblleft{}Unease; agitation\textquotedblright{}.) Say it once more.
\end{tcolorbox}
